\title{LongRoPE: Extending LLM Context Window Beyond 2 Million Tokens}

\begin{abstract}
Expanding the context window is highly sought after in large language models (LLMs). Nonetheless, present efforts to extend context windows face constraints, chiefly due to the expensive nature of fine-tuning, limited availability of sufficiently lengthy text data, and destabilizing effects caused by novel token positions. As a result, state-of-the-art extended context windows generally remain capped at approximately 128k tokens.

This work presents LongRoPE, a novel approach that, for the first time, expands the context window of pre-trained LLMs to an extraordinary \textbf{2048k} tokens. Remarkably, this is accomplished with only up to 1,000 fine-tuning steps at training lengths capped at 256k, all while preserving the original short context window performance. The method relies on three primary innovations: (i) we uncover and harness two varieties of non-uniformities in positional interpolation via an efficient search procedure, allowing for superior fine-tuning initialization and facilitating an eightfold context window extension even without fine-tuning; (ii) we propose a stepwise extension procedure that first adapts a model to a 256k context length, followed by a second stage of positional interpolation on the fine-tuned model to extend the window to 2048k tokens; (iii) we further calibrate LongRoPE at 8k tokens to restore and maintain the initial short context window effectiveness. Comprehensive empirical evaluations using LLaMA2 and Mistral on diverse benchmarks validate the utility of the proposed approach. LongRoPE-based models preserve the native model architecture with only minor alterations to the positional embedding component, and largely remain compatible with prior optimizations. The implementation will be released at \texttt{https://github.com/microsoft/LongRoPE}

\section{Introduction}

Although Large Language Models (LLMs) have achieved impressive results across a variety of tasks~\cite{gpt4,llama2}, they are typically constrained by a finite context window—for example, LLaMA2 supports up to 4096 tokens~\cite{llama2}. When required to process input sequences exceeding this window, the model's effectiveness deteriorates, largely because it lacks exposure to such extensive contexts during training. This limitation creates significant hurdles in applications such as in-context learning with many exemplars~\cite{Huang2023FewerIM}, as well as in the functioning of LLM agents~\cite{park2023generative,madaan2023selfrefine}.

Recent studies demonstrate that by fine-tuning pre-trained LLMs on lengthier textual inputs, their context window can be increased to approximately 128k tokens~\cite{longlora,pi,yarn,zhang2024soaring,liu2023scaling}. However, scaling the context window beyond this presents three primary challenges. Firstly, integrating previously unseen position indices brings in numerous anomalous values, resulting in out-of-distribution complications and hindering the stability of fine-tuning~\cite{pi}. This difficulty is pronounced when expanding from 4k to $>$1000k tokens, as this process introduces over 90\% new positions. Secondly, effective fine-tuning generally necessitates access to long documents with lengths comparable to the extended context window. Yet, existing datasets are scarce in texts surpassing 1000k tokens, and training with these lengthy sequences is computationally demanding, requiring substantial time and extensive GPU resources. Lastly, stretching the context window to extreme lengths causes the attention mechanism to become diluted, distributing focus over many token positions, which subsequently diminishes performance on tasks involving the original shorter context lengths~\cite{pi}.

To address the initial challenge, one strategy involves interpolating the RoPE positional embedding~\cite{rope,pi}, which rescales new position indices so they fit within the pretrained domain, as illustrated in Fig.\ref{fig:overview}. Position Interpolation (PI)~\cite{pi} employs a linear method to interpolate the rotary angles of RoPE based on the extension ratio. In contrast, NTK~\cite{ntk,dynamicntk} proposes non-uniform interpolation and extrapolation techniques across the different RoPE dimensions. YaRN~\cite{yarn} divides the RoPE dimensions into three distinct groups determined by frequency and implements extrapolation, NTK, and linear interpolation accordingly. Despite these efforts, positional embedding in the Transformer framework reveals a \emph{complex, heterogeneous information entropy}. The subtle irregularities in information distribution are overlooked by current methods, resulting in loss of information and thus imposing limitations on the attainable context window size.

Section~\ref{sec:analysis} presents two main empirical observations: \textbf{(1)} Successful positional interpolation must account for two distinct kinds of non-uniformity: variability across RoPE dimensions and across token positions. Less interpolation is advantageous for both lower RoPE dimensions and early token positions, though the optimal configuration varies according to the intended sequence extension length.
  \textbf{(2)} Incorporating these sources of non-uniformity into the design of positional interpolation allows more of the original RoPE structure to be preserved, especially for key dimensions and token positions. This approach effectively reduces the degradation typically introduced by interpolation, thus offering better starting points for subsequent fine-tuning. Additionally, it enables up to an 8$\times$ extension even in scenarios where fine-tuning is not performed.

Building on these results, we propose \textbf{LongRoPE}, a robust technique capable of expanding the context window in LLMs to over 2 \emph{million} tokens. The approach consists of three main components. To begin with, {LongRoPE} leverages the multidimensional non-uniform characteristics inherent in positional interpolation, determining suitable rescale factors for the rotational angles within RoPE for each dimension according to the locations of the tokens. Given that the search space for rescale factors grows exponentially relative to the desired window extension, {LongRoPE} employs an evolutionary search algorithm enhanced by two optimization strategies to improve search efficacy. Fig.~\ref{fig:overview} presents an illustration of the resulting rescaled RoPE discovered by the algorithm.

Subsequently, {LongRoPE} employs a progressive, resource-efficient extension method to reach a 2048k context window, circumventing the necessity for direct fine-tuning on ultra-long texts, which are both rare and difficult to procure. The approach initiates by identifying a suitable 256k context length on the pre-trained LLM, after which fine-tuning is carried out at this length. Next, by leveraging the capability of our non-uniform positional interpolation—which permits an 8$\times$ increase in context window without further fine-tuning—a second stage of searching is performed to determine optimal RoPE rescale factors for the already fine-tuned, length-extended LLM. Through this process, the 2048k context window is ultimately attained for both LLaMA2 and Mistral~\cite{mistral}.

In order to reduce the drop in performance on the native (shorter) context windows, {LongRoPE} modifies the RoPE rescale factors within the expanded LLM. Following the same approach used when increasing the context window size from 256k to 2048k, we also decrease the context window to 4k and 8k on the LLM that was fine-tuned for 256k, applying our search algorithm to minimize positional interpolation. When running inference, if the sequence length falls below 8k, RoPE is updated using the rescale factors found by our search.

Comprehensive evaluations conducted on multiple LLMs and a range of long-context tasks validate the superiority of our approach. Our results indicate that LongRoPE sustains low perplexity across evaluation lengths spanning from 4k to 2048k, surpasses 90\% in passkey retrieval accuracy, and attains performance on par with baselines when assessed on established benchmarks created within a 4096-token context. Furthermore, LongRoPE is compatible with all LLMs utilizing RoPE embedding. We plan to make our implementation and the LongRoPE-2048k model publicly available.

\section{Non-uniformity in Positional Interpolation}
\label{sec:analysis}

\subsection{Preliminary}

Transformer models require explicit positional information, often in the form of position embedding, to represent the order of input tokens. Our work focuses on the RoPE~\cite{rope} position embedding, which is widely used in recent  LLMs. For a token at position index $n$, its corresponding RoPE encoding can be simplified as follows:

\begin{equation}
		\fontsize{7.8}{7.8} \selectfont
	\label{eq:rope}
	\left[ cos(n\theta_0), sin(n\theta_0), cos(n\theta_1), \cdot\cdot\cdot, cos(n\theta_{d/2-1}), sin(n\theta_{d/2-1}) \right]   
\end{equation}

where $d$ denotes the embedding dimension, $n\theta_i$ is the rotation angle assigned to the token at position $n$, and $\theta_i=\theta^{-2i/d}$ specifies the corresponding rotational frequencies. In RoPE, the typical value chosen for $\theta$ is 10000.

\textbf{\emph{Context window extension ratio $s$} and positional interpolation}. Here, $s$ is defined as the quotient of the extended context length $L'$ divided by the original length $L$: $s=\frac{L'}{L}$.

To extend the context window from $L$ to $L'$, current positional interpolation methods suggest downscaling rotation frequencies $\theta_i$ based on extension ratio $s$. Let $\beta=\theta^{2/d}$, and $\lambda$ denote the actual rescale factor related to $s$, we   unify these positional interpolation methods as follows:

\begin{equation}	
	\fontsize{7.5}{7.5} \selectfont
	\label{eq:unified_rope}
	\begin{aligned}
		\left[cos\left(\frac{n}{\lambda(\beta)^0}\right), sin \left(\frac{n}{\lambda(\beta)^0}\right), cos\left(\frac{n}{\lambda(\beta)^1}\right), \cdot\cdot\cdot,  sin \left(\frac{n}{\lambda(\beta)^{d/2-1}}\right) \right]   
	\end{aligned}
\end{equation}

\textbf{Linear positional interpolation (PI)}. The method known as PI~\cite{pi} involves applying a linear interpolation to position indices, staying within the originally pre-trained sequence length. Given a desired sequence length extension by a ratio $s$, PI uniformly decreases the rotation angles for all positions by a scaling factor $\lambda=s$ throughout every RoPE dimension. Yet, this process compresses positional representations, causing position encodings for nearby tokens to be overly concentrated. This compression diminishes the model’s capability to differentiate tokens with adjacent positions, resulting in poor performance of PI as the extension ratio increases.

\textbf{NTK-based interpolation and extrapolation}. Prior works~\cite{ntk,dynamicntk} analyze RoPE through the lens of information encoding, employing the Neural Tangent Kernel (NTK) framework~\cite{jacot2018neural,tancik2020fourier}. To address the problem of positional crowding inherent in PI, they propose distributing the effects of interpolation across the different RoPE dimensions. In this approach, lower (high frequency) dimensions are scaled less while higher (low frequency) dimensions receive greater scaling. This enables both interpolation and extrapolation of position representations, following $\lambda=s^{i}$. The enhanced dynamic NTK method~\cite{dynamicntk} further refines this strategy by adapting the extension ratio per position, according to the current sequence length. In contrast to PI, which requires additional fine-tuning, NTK-aware approaches facilitate the extension of context windows without fine-tuning; however, these methods typically support a maximum extension ratio of about 4$\times$.

\textbf{YaRN}~\cite{yarn} presents a notable advance in the effectiveness of positional interpolation. The method separates the RoPE dimensions according to frequency, creating three distinct groups, each assigned a unique interpolation method. Dimensions with high frequencies are extrapolated ($\lambda$=1), low-frequency dimensions are subject to linear interpolation (PI), and those of intermediate frequency utilize the NTK approach. A central aspect of YaRN is this frequency-based partitioning of RoPE dimensions; however, the current division scheme is guided by human-driven empirical tuning, which may not yield optimal outcomes for emerging LLM architectures.

\subsection{Study on Non-uniform Positional Interpolation}

Drawing from the approaches of NTK and YaRN, which leverage non-linear characteristics—particularly by addressing frequency variation within RoPE dimensions to enhance both interpolation and extrapolation—we observe that most existing non-linear mechanisms depend predominantly on manually crafted heuristics. This observation leads us to pose two important inquiries: (1) Does the existing positional interpolation represent the best possible method? (2) Might there be alternative non-linearities that have yet to be investigated?

In order to address these inquiries, we employ evolution search (refer to Sec.~\ref{sec:method}) to identify improved non-uniform positional interpolation strategies for LLaMA2-7B. Perplexity serves as the optimization metric, based on five randomly selected excerpts from the PG19~\cite{pg19} validation dataset. Our experimental results highlight several important observations.

\textbf{Finding 1}: \textit{RoPE dimensions exhibit substantial non-uniformities, which are not effectively handled by current positional interpolation methods.}

We optimize $\lambda$ for each RoPE dimension as specified in Eq.~\ref{eq:unified_rope}. Table~\ref{tbl:exp1} presents a comparison of the perplexity scores achieved by LLaMA2-7B across various approaches on the PG19 and Proof-pile~\cite{proof-pile} test datasets, all evaluated without fine-tuning. The results from our search indicate notable gains, implying that both the conventional linear (PI) interpolation and non-uniform strategies (Dynamic-NTK and YaRN) do not yield the best results. Interestingly, on the PG19 test set, YaRN performs worse than both PI and NTK, primarily because it fails to achieve the designated context window for models that have not undergone fine-tuning. For instance, with a maximum context size of 8k, the perplexity observed with YaRN sharply increases beyond 7k.

Our investigation reveals that the scaling coefficients $\lambda$ in Eq.~\ref{eq:unified_rope} are variable across positions, in contrast to the constant scale $s$ used in PI, the formulaic approach of NTK, and YaRN's calculation within defined groups. This variability in scaling factors enhances LLaMA2's ability to model language (as measured by perplexity) for context lengths of 8k and 16k, even when no fine-tuning is applied. The core advantage stems from the positional embedding maintaining the integrity of RoPE—most notably in crucial dimensions—which in turn lessens the model's challenges in differentiating token positions that are spatially proximate.

\textbf{Finding 2}: \textit{RoPE for the initial tokens in the input sequence should be extrapolated with less interpolation.} 

We propose that the RoPE mechanism should perform minimal interpolation for the first $\hat{n}$ tokens of each input sequence, as these tokens generally attract higher attention scores and are pivotal within the attention mechanism, in line with findings from Streaming LLM~\cite{streamingllm} and LM-Infinite~\cite{han2023lminfinite}. To test this idea, we expand the context window to 8k and 16k tokens using PI and NTK approaches, ensuring that the initial $\hat n$ tokens (ranging from 0, 2, ..., 256) bypass interpolation. Notably, setting $\hat n$=0 defaults to the original PI and NTK configurations. Results presented in Table~\ref{tbl:tokenposition} reveal two major insights: \textbf{(1)} Excluding the starting tokens from position interpolation enhances overall model performance; \textbf{(2)} The optimal value for $\hat n$ varies according to the desired length of the context extension.

\textbf{Finding 3}: \textit{Non-uniform positional interpolation effectively extends LLM context window in both fine-tuning and non-fine-tuning settings.} 

Although our experiments demonstrated that the non-uniform position interpolation we identified offers notable improvements in extension capabilities at 8k and 16k context lengths without the need for additional training, extending to even longer context windows necessitates fine-tuning. Consequently, we apply fine-tuning to LLaMA2-7B using our discovered RoPE configuration for sequences up to 64k tokens (see the Appendix for experimental details). According to the results presented in Table~\ref{tbl:finetune}, our approach yields considerable performance gains over PI and YaRN, both in the original model and following fine-tuning. These advantages stem from the efficient exploitation of non-uniform positional interpolation, which reduces the loss of information and supplies superior initial model parameters for subsequent fine-tuning.

\textbf{Summary}. This work identifies two types of non-uniformity: discrepancies in RoPE dimensions and variations across token positions. Harnessing these non-uniformities in the design of positional interpolation leads to substantial enhancements in the ability of LLMs to handle extended context lengths.

\section{LongRoPE}

\label{sec:method}
Drawing on these insights, we propose LongRoPE, a framework that initially incorporates a streamlined search algorithm to capitalize on the two distinct non-uniformities. This algorithm is then leveraged to expand the LLM context window past 2 million tokens.

\subsection{Problem Formulation}

These two forms of non-uniformity expand the solution space significantly and complicate the optimization process. To mitigate these challenges, we reformulate the multidimensional non-uniform position interpolation optimization task as a search problem.

For a LLM targeting a  context window size of $L'$ and lengthy input documents $\mathbf{X}$, where each $\mathbf{x}\in\mathbf{X}$ surpasses $L'$ in token length, we denote the original rotary angle of the $i^{th}$ dimension in RoPE embedding at token position $n$ as  $\frac{n}{\beta^i}$. The optimization problem is then formulated as follows:

\begin{equation}
\begin{gathered}
		\label{eq:problem}

\underset {\mathbf{x}\in\mathbf{X}; \, |\mathbf{x}|\ge L'}{ \operatorname {arg\,min} } \,  \mathcal{L}\, \left( \text{LLM} (\text{RoPE}, \mathbf{X})\right), \text{ such that }
\\
\scalebox{0.93}{
$\underset {\begin{subarray}{l} i=0,\ldots,\frac{d}{2}-1;\\ n\in[0, |\mathbf{x}|);\end{subarray}}{\text{RoPE}_(n)} = {\left[\cdots, \cos\left(\mathbb{I}(\hat{\lambda}_{i}, \hat{n}) \frac{n}{\beta^i}\right), \sin\left(\mathbb{I}(\hat{\lambda}_{i}, \hat{n}) \frac{n}{\beta^i}\right), \cdots\right]}$}\\
\text{where} \; \mathbb{I}(\hat{\lambda}_{i},\hat{n})= 
\begin{cases}
1 &\text{if } n<\hat{n} \\
\frac{1}{\lambda_i} &\text{if } n\geq\hat{n}
\end{cases}
	\end{gathered}
\end{equation}

We define a collection of scaling coefficients, $\mathbb{I}(\hat{\lambda}_{i},\hat{n})$, in order to account for the two types of non-uniformity present. The symbols $\hat{\lambda_i}$ and $\hat{n}$ are used to represent the non-uniform scaling among RoPE dimensions and across token positions, respectively. More precisely, $\mathbb{I}(\hat{\lambda}_{i},\hat{n})$ serves to adjust the rotation angle for the $i^{th}$ RoPE dimension, with $\hat{\lambda}_i$ specifying the dimension-specific scaling and $\hat{n}$ indicating a positional threshold for tokens. For the first $\hat{n}$-$1$ tokens, the scaling coefficient $\hat{\lambda}_i$ does not alter the rotation, thus the standard RoPE angle $\frac{n}{\beta^i}$ is preserved. Once token positions reach $n\ge\hat{n}$, the scaling factor is then utilized.

Assuming a desired context window length $L'$, the goal is to determine the best rescale factors ($\mathbb{I}(\hat{\lambda}_{0},\hat{n})$, $\mathbb{I}(\hat{\lambda}_{1},\hat{n})$, ..., $\mathbb{I}(\hat{\lambda}_{i},\hat{n})$, ...) across all RoPE dimensions, ranging from the $1^{st}$ to the $d^{th}$. By applying these rescaled values to the target $\text{LLM}$'s $\text{RoPE}$ mechanism, the model should attain the lowest possible next token prediction loss, denoted as $\mathcal{L}$ (corresponding to perplexity), on input sequences $\mathbf{X}$ with token length $L'$.

\subsection{Searching the Non-uniform Position Interpolation}

To address the issue presented in Eq.~\ref{eq:problem}, we propose a straightforward but powerful approach that seeks optimal RoPE rescale factors, aiming to capitalize on the inherent multidimensional disparities present within position embeddings.

\textbf{Search space}. We construct an extensive search space that accounts for the two types of non-uniformities. This is detailed in Table~\ref{tbl:searchspace}. In particular, the configuration permits independent optimization of rescale factors for each RoPE embedding dimension. To streamline the structure of the search space, we opt to search for $\lambda_i$ and $\hat{n}$ directly, rather than exploring the indicator function $\mathbb{I}(\hat{\lambda}_{i},\hat{n})$, where $\hat{\lambda}_{i}=1/\lambda_i$. As depicted in Table~\ref{tbl:searchspace}, the parameter $\lambda_i$ is selected from a lower bound of 1.0 (representing straightforward extrapolation) up to an upper bound of $s\times1.25$ (enabling greater interpolation compared to PI) with increments of 0.01. Here, $s$ corresponds to the target ratio for extending the context window.

The parameter $\hat{n}$ determines how many of the first token positions are preserved without applying positional interpolation, instead relying on the original RoPE embedding. In practice, $\hat{n}$ is selected from the set \{0, 1, 2, 4, 8, 12, 16, 20, 24, 28, 32, 64, 128, 256\}. If $\hat{n}$ is set to 0, rescale factors obtained from the search are utilized for all token positions.

\textbf{Evolution-based search}. The search space outlined in Table~\ref{tbl:searchspace} encompasses a vast array of positional interpolation options, making thorough and efficient search a daunting task. For instance, when considering an extension factor of $s=4\times$, the total number of possible configurations can reach $400^{128/2}\times14$=4$\times10^{167}$. As the extension ratio increases, this search space grows at an exponential rate. To manage this complexity, we employ evolution search~\cite{guo2020single}, incorporating two specialized optimization strategies to significantly enhance search efficiency. The entirety of the search algorithm is detailed in Algorithm~\ref{alg:search}.

\textit{Optimized initial population generation}. Rather than starting with a purely random set of $P$ rescale factors for the population, we ensure that the initial group includes the three specific RoPE rescale factors associated with PI, NTK, and YaRN as dedicated members. The other $P$-3 members of the population are produced by applying random mutations to these three reference rescale factors, each with a probability $p$.

\textit{Monotonically non-decreasing constraint}. Once the initial population is created, each individual’s LLM perplexity is evaluated. We assign the relevant RoPE rescale factors within the target LLM and assess perplexity using the input $\mathbf{X}$. The best-performing $k$ individuals are selected as parents for subsequent generations. Despite this, random mutations and crossover in a large search space may frequently yield low-quality candidates, resulting in redundant perplexity evaluations. This inefficiency becomes more pronounced as $L'$ increases, since every perplexity assessment incurs substantial computational cost.

To tackle this issue, we enforce a monotonic non-decreasing constraint on the RoPE rescaling factors, such that $\lambda_i\le \lambda_{i+1}$ holds true for all $i$. RoPE configurations that conform to this requirement are used in evaluating LLM perplexity, which streamlines and greatly cuts down on the computational burden associated with the search process. More precisely, we mandate that the sequence $\lambda_i$ grows monotonically as the dimension $i$ (from $0$ to $63$) increases. The rationale for this monotonicity by dimension arises from NTK theory~\cite{jacot2018neural,tancik2020fourier,ntk}: lower-indexed dimensions, corresponding to higher-frequency components, necessitate minimal interpolation (hence, smaller $\lambda_i$), whereas higher-indexed dimensions, linked to lower frequencies, tolerate greater interpolation (larger $\lambda_i$).

\begin{algorithm}[t]
	\small
	\caption{The search algorithm}
	\textbf{Input:} target LLM, input samples $\mathbf{X}$, population size $P$, mutation size $N_1$, crossover size $N_2$, maximum iterations $\mathcal{T}$, mutation probability $p$\\

	\begin{algorithmic}[1]
		\label{alg:search}
		\STATE Initialize Top-k as empty;	
		\STATE Set $\text{P}_0$ by invoking \textit{Initialize\_population\_with\_optimization}($P$, $\mathbf{X}$, $p$); 
		\FOR{$i$ from $1$ to $\mathcal{T}$}
		\STATE Execute \textit{Compute\_perplexity} (LLM, $\text{P}_{i-1}$, $\mathbf{X}$);
		\STATE Update Top-k using \textit{Update\_Topk} (Top-k, $\text{P}_{i-1}$);
		\STATE Generate $\text{P}_{mutation}$ via \textit{Mutation\_with\_mono\_constraint} (Top-k, $N_1$, $p$); 
		\STATE Obtain $\text{P}_{crossover}$ through \textit{Crossover\_with\_mono\_constraint} (Top-k, $N_2$);
		\STATE Formulate $\text{P}_i$ by merging $\text{P}_{mutation}$, $\text{P}_{crossover}$, and Top-k;
		\ENDFOR
		\STATE Output the member from Top-k with minimal perplexity;
	\end{algorithmic}
\end{algorithm}

\textbf{8$\times$ extension without fine-tuning}. Through the use of evolutionary search, our approach successfully determines optimal non-uniform RoPE rescale factors, which help maintain crucial dimensions and position representations, thereby reducing information degradation due to interpolation. As illustrated in Fig.\ref{fig:non-ft}, this strategy enables the LLaMA2 model's context window to be expanded from 4k to 32k tokens without the need for fine-tuning. By comparison, prior approaches such as PI, as well as non-uniform NTK and YaRN, experience significant increases in perplexity after only a 2$\times$ extension.

\subsection{Extending LLM Context Window to 2048K}
\label{sec:extendto2048k}

\textbf{Progressive extension to 2048k}. In this section, we present our strategy for expanding the context window of pre-trained LLMs well beyond the standard 4k, reaching up to and exceeding 2048k. Our approach leverages non-uniform positional interpolation, allowing an 8$\times$ increase in context length without necessitating fine-tuning. However, when aiming for more substantial extensions (for instance, up to 512$\times$), fine-tuning becomes indispensable. A common technique involves identifying optimal RoPE rescaling factors at the desired 2048k sequence length followed by fine-tuning. Nonetheless, this process is hampered by the considerable computational demands it entails. Furthermore, in our experiments, we have observed that fine-tuning LLMs under such large extension ratios is nontrivial and often proves challenging (refer to Appendix for further details).

Fortunately, LongRoPE demonstrates robust performance with both pre-trained and further fine-tuned extended LLMs. To this end, we propose a streamlined and incremental approach that reaches the intended 2048k context window using only 1k fine-tuning iterations, all conducted within a 256k training sequence length.

$\Diamond$ \textit{LongRoPE search for expanding pre-trained LLMs to 256k}. Using LLaMA2 as a representative case, we explore search strategies for enlarging the context window to both 128k and 256k tokens. At this point, the context length increases by factors of 32$\times$ and 64$\times$, correspondingly.

$\Diamond$ \textit{Fine-tuning to 256k}. Following pre-training, we adapt the LLM to support a context window of 256k. The process begins with fine-tuning LLaMA2 for 400 steps employing the RoPE rescaled factors appropriate for 128k. After this stage, we update the RoPE rescaled factors to those for 256k on the obtained checkpoint and perform an extra 600 fine-tuning steps. This sequential strategy demonstrates greater efficiency compared to immediately fine-tuning at 256k.

$\Diamond$ \textit{Scaling fine-tuned extended LLM to 2048k via LongRoPE search}. In the final stage, an additional search process is applied to the LLM previously fine-tuned at 256k sequence length. This approach enables us to reach a vast context window of 2048k, requiring no extra rounds of fine-tuning. As a result, the total extension factor is $512\times$.

\textbf{Recovery for shorter context windows}. Following the expansion to a substantial 2048k context window, we observe that model performance degrades within the initial context range. This outcome aligns with established limitations of positional interpolation~\cite{pi}, which compels the position embeddings in the lower context dimensions to be compressed into a more confined space, thereby diminishing model effectiveness. When the extension ratio reaches 512$\times$, positions within the original 4k context are densely packed, exacerbating this effect.

To address this issue, an additional evolutionary search is conducted on the expanded LLM, specifically to tune RoPE rescale factors at shorter context lengths such as 4k and 8k. The upper bound for the searched $\lambda$ is lowered in this process, reflecting the diminished need for positional interpolation in these shorter scenarios. At inference time, the LLM adaptively modifies the relevant RoPE rescale factors based on the context length in use.

\section{LongRoPE}

\label{sec:method}
Building upon these insights, we propose LongRoPE, which incorporates a novel search algorithm designed to effectively utilize both identified non-uniformities. This approach subsequently enables the extension of the LLM context window to exceed 2 million tokens.

\subsection{Problem Formulation}

These two types of non-uniformity expand the solution space significantly, resulting in increased complexity during optimization. To mitigate this, we reformulate the multidimensional non-uniform position interpolation optimization challenge as a search problem.

For a LLM targeting a  context window size of $L'$ and lengthy input documents $\mathbf{X}$, where each $\mathbf{x}\in\mathbf{X}$ surpasses $L'$ in token length, we denote the original rotary angle of the $i^{th}$ dimension in RoPE embedding at token position $n$ as  $\frac{n}{\beta^i}$. The optimization problem is then formulated as follows:

\begin{equation}
\begin{gathered}
		\label{eq:problem}

\underset {\mathbf{x}\in\mathbf{X}; \, |\mathbf{x}|\ge L'}{ \operatorname {arg\,min} } \,  \mathcal{L}\, \left( \text{LLM} (\text{RoPE}, \mathbf{X})\right),\quad \text{where} 
\\
\scalebox{0.93}{
$\underset {\begin{subarray}{l} i=0,\ldots,\frac{d}{2}-1;\\ n\in[0,|\mathbf{x}|);\end{subarray}}{\text{RoPE}_(n)}= {\left[\ldots,\cos\left(\mathbb{I}(\hat{\lambda}_{i}, \hat{n})\times\frac{n}{\beta^i}\right),  \sin\left(\mathbb{I}(\hat{\lambda}_{i}, \hat{n})\times\frac{n}{\beta^i}\right),\ldots\right] }$}\\
	\text{where} \quad \mathbb{I}(\hat{\lambda}_{i},\hat{n})=\begin{cases}
	1&\text{\; if $n<\hat{n}$}\\
{\frac{1}{\lambda}_{i}}&\text{\; if $n\geq\hat{n}$}
\end{cases}
	\end{gathered}
\end{equation}

In this approach, we define a set of rescale factors, $\mathbb{I}(\hat{\lambda}_{i},\hat{n})$, to address both varieties of non-uniformity. Here, $\hat{\lambda}_{i}$ and $\hat{n}$ represent the irregularities in RoPE dimensions and token positions, respectively. The function $\mathbb{I}(\hat{\lambda}_{i},\hat{n})$ is employed to adjust the rotation angle corresponding to the $i^{th}$ RoPE dimension, with $\hat{\lambda}_{i}$ serving as the scaling coefficient and $\hat{n}$ as the threshold for token positions. For the first $\hat{n}-1$ token locations, the rescale coefficient $\hat{\lambda}_{i}$ remains inactive, utilizing the original RoPE rotary angle $\frac{n}{\beta^i}$. For all tokens at positions $n \geq \hat{n}$, the rescale factor becomes effective.

For a given target context window length $L'$, the aim is to determine the set of optimal rescaling factors—namely, $\mathbb{I}(\hat{\lambda}_{0},\hat{n})$, $\mathbb{I}(\hat{\lambda}_{1},\hat{n})$, ..., $\mathbb{I}(\hat{\lambda}_{i},\hat{n})$, up to the $d^{th}$ RoPE dimension. Implementing these scaling factors in the $\text{RoPE}$ mechanism of the target $\text{LLM}$ should yield the lowest possible loss in next token prediction, $\mathcal{L}$ (i.e., perplexity), across input datasets $\mathbf{X}$ that have a sequence length of $L'$.

\subsection{Searching the Non-uniform Position Interpolation}

To address the challenge posed in Eq.~\ref{eq:problem}, we propose a straightforward but powerful approach that identifies the best RoPE rescale factors, allowing us to leverage the varying multi-dimensional characteristics of position embeddings to their fullest extent.

\textbf{Search space}. We construct an extensive search space that incorporates both types of non-uniformity. Table~\ref{tbl:searchspace} provides an overview of this search space. In particular, the framework enables tuning of individual rescale factors for each dimension within RoPE embedding. To facilitate straightforward search space construction, the variables $\lambda_i$ and $\hat{n}$ are directly optimized, instead of searching over $\mathbb{I}(\hat{\lambda}_{i},\hat{n})$, where $\hat{\lambda}_{i}$ is defined as $1/\lambda_i$. As detailed in Table~\ref{tbl:searchspace}, $\lambda_i$ is selected within the interval starting from 1.0 (representing direct extrapolation) up to $s\times1.25$ (representing a broader range of interpolation compared to PI), with increments of 0.01. Here, $s$ denotes the extension ratio of the intended context window.

The parameter $\hat{n}$ specifies how many of the starting token positions maintain their original RoPE embeddings, without any position interpolation. Based on empirical observations, we conduct a search over $\hat{n}$ values in the set \{0, 1, 2, 4, 8, 12, 16, 20, 24, 28, 32, 64, 128, 256\}. If $\hat{n}=0$, this indicates that all token positions apply the optimized rescale factors.

\textbf{Evolution-based search}. The search space detailed in Table~\ref{tbl:searchspace} encompasses a vast array of positional interpolation methods, making thorough exploration highly challenging. For instance, when using a $s=4\times$ extension, the number of possible configurations becomes $400^{128/2}\times14$=4$\times10^{167}$, resulting in an astronomical search space. As the extension ratio increases, the size of this space grows exponentially. To mitigate these challenges, we employ evolution search~\cite{guo2020single} and introduce two novel optimization strategies that substantially improve the efficiency of the search process. The main workflow is presented in Algorithm~\ref{alg:search}.

\textit{Improved initialization of the population}. In lieu of randomly generating an entire population of $P$ rescale factors, our approach includes the three specific RoPE rescale factors associated with PI, NTK, and YaRN as predefined individuals within the initial population. The other $P-3$ individuals are produced by applying random mutations to these three rescale factors, each with probability $p$.

\textit{Monotonically non-decreasing constraint}. Once the initial population is created, we evaluate the LLM perplexity for each member. This involves implementing the specific RoPE rescale factors within the target LLM and calculating the perplexity for input $\mathbf{X}$. The best-performing $k$ individuals are then selected as parents for the evolutionary process. Despite this, the immense search space means that uninformed mutation and crossover steps frequently result in the selection of suboptimal candidates, prompting superfluous perplexity evaluations. This inefficiency becomes especially problematic as $L'$ increases, since each perplexity measurement requires considerable inference time.

To mitigate this issue, we enforce a monotonicity restriction such that the rescaled RoPE factors are non-decreasing: $\lambda_i\le \lambda_{i+1}$. Only RoPE configurations that adhere to this constraint are utilized during the LLM perplexity assessment, substantially limiting the computational overhead associated with the search process. Concretely, we mandate that $\lambda_i$ follows a monotonically increasing pattern across the RoPE dimensions, with $i$ ranging from 0 to 63. This design principle draws from NTK theory~\cite{jacot2018neural,tancik2020fourier,ntk}, which posits that dimensions associated with higher frequencies require less interpolation (smaller $\lambda_i$), whereas dimensions corresponding to lower frequencies can tolerate greater interpolation (larger $\lambda_i$).

\begin{algorithm}[t]
	\small
	\caption{The search algorithm}
	\textbf{Input:} target LLM, input samples $\mathbf{X}$, population size $P$, mutation size $N_1$, crossover size $N_2$, max iterations $\mathcal{T}$, mutation probability $p$\\

	\begin{algorithmic}[1]
		\label{alg:search}
		\STATE Top-k $\leftarrow$ $\phi$;	
		\STATE $\text{P}_0$ $\leftarrow$ \textit{Initialize\_population\_with\_optimization}($P$, $\mathbf{X}$, $p$); 
		\FOR{$i = 1$ to $\mathcal{T}$}
		\STATE \textit{Compute\_perplexity}(LLM, $\text{P}_{i-1}$, $\mathbf{X}$);
		\STATE Top-k $\leftarrow$ \textit{Update\_Topk}(Top-k, $\text{P}_{i-1}$);
		\STATE $\text{P}_{mutation}$ $\leftarrow$ \textit{Mutation\_with\_mono\_constraint}(Top-k, $N_1$, $p$); 
		\STATE $\text{P}_{crossover}$ $\leftarrow$ \textit{Crossover\_with\_mono\_constraint}(Top-k, $N_2$);
		\STATE $\text{P}_i$ $\leftarrow$ $\text{P}_{mutation}$ $\cup$ $\text{P}_{crossover}$ $\cup$ Top-k;
		\ENDFOR
		\STATE Output the individual from Top-k that achieves minimum perplexity;
	\end{algorithmic}
\end{algorithm}

\textbf{8$\times$ extension without fine-tuning}. Through the use of evolutionary search, our approach efficiently uncovers optimal non-uniform RoPE rescale factors that maintain crucial positions and dimensions, thereby reducing the information loss resulting from interpolation. As shown in Fig.\ref{fig:non-ft}, this enables our method to expand the context length of LLaMA2 from 4k up to 32k tokens without any requirement for fine-tuning. In comparison, alternative techniques like PI, as well as non-uniform NTK and YaRN, exhibit a sharp increase in perplexity once the context window surpasses a 2$\times$ extension.

\subsection{Extending LLM Context Window to 2048K}
\label{sec:extendto2048k}

\textbf{Progressive extension to 2048k}. In this section, we present an approach to expand the context length of pre-trained LLMs beyond the conventional 4k, reaching up to 2048k. Our experiments reveal that employing non-uniform positional interpolation permits an 8$\times$ increase in context length without the need for additional fine-tuning. When aiming for even greater expansions, such as up to 512$\times$, fine-tuning becomes indispensable. A possible strategy involves searching for RoPE rescaling parameters compatible with the target 2048k context and proceeding with fine-tuning thereafter. Nonetheless, this process is hindered by the substantial computational resources required for training. In addition, our observations indicate that effectively fine-tuning LLMs for such significant extensions is particularly challenging (refer to Appendix).

Fortunately, LongRoPE demonstrates efficacy in both the pre-trained and fine-tuned forms of the extended LLM. Consequently, we propose an efficient and incremental approach that reaches the desired 2048k length utilizing only 1k fine-tuning iterations, all conducted at a training length capped at 256k.

$\Diamond$ \textit{LongRoPE-based context window extension of pre-trained LLM to 256k}. Using LLaMA2 as a case study, we search for configurations that support context window sizes of 128k and 256k tokens. In these scenarios, the model’s receptive field is enlarged by factors of 32$\times$ and 64$\times$, respectively.

$\Diamond$ \textit{Fine-tuning to 256k}. Following pre-training, we adapt the LLM to handle a context window of 256k through a staged fine-tuning process. Initially, LLaMA2 is fine-tuned for 400 steps utilizing RoPE rescaled factors corresponding to 128k. After this phase, we update the checkpoint by switching to RoPE rescaled factors for 256k and continue with 600 additional fine-tuning steps. This two-stage approach is more effective than performing all fine-tuning directly at the 256k context size.

$\Diamond$ \textit{Expanding the fine-tuned extended LLM to 2048k using LongRoPE search}. In the last stage, we carry out a follow-up search procedure applied to the LLM that has already been fine-tuned for a context length of 256k. This approach successfully pushes the maximum context window to 2048k, all without additional rounds of fine-tuning. Consequently, the overall extension factor reaches $512\times$.

\textbf{Recovery of performance in shorter context windows}. Upon expanding the context window to a notably large 2048k, we observe a decline in model performance when evaluated on the initial context range. This decline has been previously documented in the context of positional interpolation~\cite{pi}; it occurs because embeddings corresponding to positions within the original context window are compressed into a limited subspace in higher dimensions, thereby impairing the language model's capability. Specifically, with a 512$\times$ scaling of the context window, positional representations in the initial 4k window become especially densely packed.

To address this issue, we conduct an additional evolution search on the extended LLM, specifically tuning the RoPE rescale factors for shorter context lengths such as 4k and 8k tokens. For these reduced lengths, we limit the upper bound of the searched $\lambda$ since they require less positional interpolation. During inference, the LLM adaptively modifies the associated RoPE rescale factors as needed.

 \section{Experiments}

\label{sec:exp}
\subsection{Setup}
\textbf{Evaluation Tasks and models}. Our experiments employ LongRoPE with LLaMA2-7B and Mistral-7B, assessing their effectiveness based on three criteria: (1) measuring the perplexity of large language models with extended context lengths on lengthy textual inputs; (2) performance in a Passkey retrieval task, designed to test the model’s capacity to extract a target key from substantial distractor text; and (3) results on widely-used LLM benchmarks with a standard context window of 4096 tokens.

\textbf{Fine-tuning}. For LLaMA2, we adopt a learning rate of 2e-5 with linear scheduling, maintaining a global batch size of 32. The model is initially fine-tuned for 400 steps on the Redpajama~\cite{redpajama} corpus, which is partitioned into segments of 128k tokens and marked by BOS and EOS boundaries. Subsequently, we extend training by an additional 600 steps, using the previously obtained checkpoint, in order to scale the context window to 256k tokens. The 128k token configuration utilizes 8 A100 GPUs leveraging the distributed system described in~\cite{cube}, whereas the 256k token setup is trained on 16 A100 GPUs. For Mistral, training is performed with a fixed learning rate of 1e-6 and a global batch size of 64. For both the 128k and 256k variants, we implement the YaRN~\cite{yarn} protocol, fine-tuning over 400 steps on Together Computer's Long-Data Collections~\cite{mistraldata} using sequences of length 16k, employing 4 A100 GPUs in each case.

\textbf{Search}. For target window sizes up to 256k, the parameters are set as follows: $P$=64, $N_1$=$N_2$=16, $p$=0.3, and $\mathcal{T}$=40. In each iteration, the top-32 candidates are chosen for mutation and crossover. Perplexity assessments are based on 5 randomly chosen validation samples from the PG19 set, ensuring each meets the minimum length defined by the target context. For window sizes above 512k, the sizes for the population, mutation, and crossover operations are reduced by half. In these cases, perplexity is computed over 3 randomly selected validation samples from the Pile-Books3~\cite{pile} dataset.

\textbf{Baselines}. In order to achieve a 2048k context window, we applied fine-tuning procedures to models initially trained with 128k and 256k context lengths. This produces LongRoPE-2048k (ft=128k) and LongRoPE-2048k (ft=256k), targeting LLaMA2 and Mistral, respectively. Our evaluation includes a comparison against several leading baseline methods for context window extension, with a focus on publicly available LLMs that have undergone positional interpolation techniques, including PI, NTK, and YaRN. The set of baselines encompasses Together-32k~\cite{together}, Code LLaMA~\cite{codellama}, LongLoRA-full-FT-100k~\cite{longlora}, YaRN-LLaMA, and YaRN-Mistral~\cite{yarn}.

\subsection{Main Results}

\label{sec:mainresult}
\textbf{Long sequence language modeling up to 256k tokens}. Our evaluation starts with a direct comparison to cutting-edge LLMs extended to accommodate up to 256k tokens. To highlight our approach’s versatility, we select two benchmark datasets: the test splits of Proof-pile~\cite{pg19} and PG19~\cite{pile}. We measure perplexity at different context lengths, implementing a sliding window of 256 tokens during assessment. For PG19, the evaluation encompasses all 100 documents present in the test set. For Proof-pile, adhering to the protocol established in YaRN~\cite{yarn}, we draw 10 samples at random, each spanning at least 128k tokens.

Table~\ref{tbl:proof-pile} and Table~\ref{tbl:pg19} present a comparison of perplexity values for LLaMA2 and Mistral models extended using various interpolation approaches on the Proof-pile and PG19 datasets, respectively. Two main findings emerge: \textbf{(1)} The perplexity of our extended models consistently declines as the evaluation length increases from 4k to 256k, demonstrating enhanced capability in utilizing extended contexts. \textbf{(2)} Notably, even when the context window is scaled up to 16$\times$ its original size—a scenario where preserving short-context performance is typically difficult—our LongRoPE-2048k models surpass state-of-the-art baselines within the 256k context range.

\textbf{Modeling languages over sequences surpassing 2000k}. To assess performance on exceptionally lengthy texts, the Books3~\cite{pile} collection is employed. To streamline the evaluation, 20 books longer than 2048k tokens are chosen at random, and a 256k sliding window is applied.

Table~\ref{tbl:books3} demonstrates that LongRoPE effectively expands the context window for LLaMA2-7B and Mistral-7B models to 2048k, achieving perplexity scores that are comparable to, or even outperform, the baselines at shorter context sizes ranging from 8k to 128k. Notably, there are distinct performance variations observed between the 2048k versions of LLaMA2 and Mistral. At shorter context windows, Mistral exhibits better performance than the baselines, although its perplexity rises above 7 when the context length exceeds 256k. LLaMA2's results behave as anticipated: perplexity improves with longer contexts, with slight upticks at 1024k and 2048k. Furthermore, LongRoPE-2048k achieves superior results on LLaMA2 when fine-tuned at 256k compared to 128k, which is attributed to the smaller secondary extension ratio (8$\times$ versus 16$\times$). Conversely, Mistral benefits more from fine-tuning at 128k. This discrepancy mainly arises because, for Mistral's 128k and 256k fine-tuning experiments, a 16k training length is used according to YaRN's methodology, thereby influencing Mistral's capacity to extend its context window post fine-tuning.

\textbf{Passkey retrieval}. Next, we investigate the influence of context window size on generation-based tasks. Adopting a synthetic assessment known as passkey retrieval from ~\cite{passkey}, the task requires the model to locate and extract a randomly chosen five-digit passkey embedded within a lengthy document. The specific prompt design can be found in the appendix. We conduct 10 independent runs of the passkey retrieval task, ensuring the passkey appears at random positions uniformly distributed throughout the evaluation context length.

Fig.~\ref{fig:passkey} presents a comparison of retrieval accuracy against baseline models. While the accuracy of current approaches declines sharply to zero once the context length exceeds 128k tokens, LongRoPE-LLaMA2-2048k (ft=256k) demonstrates robust performance. Specifically, it consistently achieves retrieval accuracy of at least 90\% across the span of 4k to 2048k tokens, even in the demanding scenario of extracting a passkey from millions of tokens. Meanwhile, LongRoPE-Mistral-2048k (ft=128k) sustains perfect accuracy up to 1800k tokens, only decreasing to 60\% at the 2048k mark. This outcome mirrors the results shown in Table~\ref{tbl:books3}, where the perplexity exhibits a slight increase at the longest context length.

\textbf{Evaluation on standard benchmarks in the original context window}. LongRoPE-2048k models are assessed within their native context window via the Hugging Face Open LLM Leaderboard~\cite{huggingfaceboard}, incorporating both zero-shot and few-shot protocols. Specifically, performance is measured using the 25-shot ARC-Challenge~\cite{allenai:arc}, 10-shot HellaSwag~\cite{zellers2019hellaswag}, 5-shot MMLU~\cite{mmlu}, and 0-shot TruthfulQA~\cite{lin2021truthfulqa}.

As demonstrated in Table~\ref{tbl:commonsense}, our models deliver similar performance to prior models on the original benchmark, which was intended for use with shorter context lengths, and surpass the original Mistral by 0.5\% on TruthfulQA. Although LongRoPE-LLaMA2-2048k—fine-tuned with a 256k context—exhibits a marginally higher drop in performance, its results still largely fall within acceptable limits across the majority of tasks.

\subsection{Ablation Results}

\textbf{Evaluating the second positional interpolation's effectiveness}. Through the implementation of our progressive extension framework, we employ our search algorithm to apply a second layer of non-uniform positional interpolation to the extended and fine-tuned LLMs. To assess the utility of this approach, we carry out experiments utilizing the fine-tuned LLaMA2-256k model, expanding its context length to 512k, 1024k, and 2048k via PI and YaRN. Table~\ref{tbl:secondsearch} demonstrates that our method of non-uniform positional interpolation maintains stable perplexity across various extension ratios. By comparison, both PI and YaRN exhibit a noticeable and rapid increase in perplexity as the context length grows.

\textbf{Recovery performance in short context windows.} To address diminished effectiveness at reduced context sizes, we re-tuned the RoPE factors for LongRoPE-2048k using our search method. In this process, we constrained the maximum scale factors permissible during search to promote minimal interpolation at shorter 4k and 8k contexts. Table~\ref{tbl:shortperformance} presents the perplexity results of LongRoPE-LLaMA2-2048k evaluated on Proof-pile at 4k and 8k, together with mean benchmark accuracy across LLM tasks. These findings indicate substantial gains at shorter context lengths.

\textbf{Evaluation of the two types of non-uniformity}. In this section, we examine the influence of each form of non-uniformity by performing ablation studies. We design two experimental settings: (i) for LLaMA2-7B, we expand the sequence length to 16k and 32k using three strategies—PI, adjusting only the RoPE dimension, and modifying both types of non-uniformities; (ii) for our fine-tuned 256k-length LLaMA2 model, we further stretch its context to 2048k with identical methodology. The assessment of perplexity is conducted without any additional fine-tuning. Table~\ref{tbl:searchdimension} illustrates that introducing non-uniformity into the RoPE dimension substantially decreases perplexity compared to the linear interpolation approach used by PI. Incorporating non-uniformity in token position yields notable gains at 16k and 32k lengths, but this effect is diminished at 2048k—likely attributable to the sheer length of the sequence. Simply keeping the original tokens without interpolation becomes ineffective in these extremely long contexts, which we propose to investigate further in subsequent research.

\section{Related Works}

Besides techniques relying on position interpolation, this section examines related work employing alternative strategies.

\textbf{Retrieval-based} strategies leverage external memory components to store extended contextual information, along with retrieval modules that fetch pertinent documents during inference~\cite{longllama,wang2023augmenting,borgeaud2022improving}. Such systems usually require substantial, explicit architectural changes to the underlying LLMs. By contrast, the approach proposed in this study is comparatively streamlined, involving only minor adjustments to positional embeddings. Additionally, our method accommodates a broader range of long context tasks, extending beyond retrieval to include areas such as extended document summarization and few-shot learning.

\textbf{Attention-based approaches for expanding the context window}. In addition to interpolating positional embeddings, several studies have extended the input context window length for LLMs by modifying their attention operations~\cite{han2023lminfinite, streamingllm,parallelwindow}. Central to these techniques is the development of innovative attention masks, which serve to address the rapid growth in attention computations that arises when introducing additional positions. These approaches complement positional interpolation strategies.

\textbf{Fine-tuning based approaches} investigate strategies to efficiently adapt pre-trained LLMs with modified position embeddings to accommodate longer sequence lengths. Research efforts such as Code LLaMA~\cite{codellama}, LLaMA2 Long~\cite{xiong2023effective}, and ScaledRoPE~\cite{liu2023scaling} utilize significantly larger base values for RoPE and perform fine-tuning on sequences matching the intended length. In contrast, our approach enables versatile adaptation to a range of target context lengths and is capable of scaling to lengths exceeding 2M tokens. Additionally, given the extensive GPU resources required for fine-tuning on very long contexts (e.g., beyond 128k), recent methods like LongLoRA~\cite{longlora} and PoSE~\cite{zhu2023pose} have been introduced to alleviate this computational burden. Our method is complementary and operates independently of these efficient fine-tuning techniques.

\section{Conclusion}

This study introduces LongRoPE, a technique that significantly increases the context window of LLMs to an extraordinary 2048k, all while preserving their performance within the original, shorter windows. We utilize two distinct types of non-uniformities inherent in RoPE positional embeddings, identified through an optimized evolutionary search approach. This strategy delivers dual advantages: it serves as a robust starting point for subsequent fine-tuning, and it allows for an 8$\times$ extension of the context window even in the absence of fine-tuning. Furthermore, we implement a stepwise expansion procedure leveraging LLMs fine-tuned with a 256k context length, ultimately attaining a 2048k context window size without necessitating additional fine-tuning. A comprehensive set of experiments demonstrates the effectiveness of LongRoPE. We anticipate that LongRoPE-2048k models will pave the way for novel applications requiring long contexts and foster future advancements in this domain.

\section*{Broader Impacts} The research discussed in this paper aims to further the progress of Machine Learning as a discipline. While our findings may have various implications for society, we do not believe that any particular aspect requires distinct emphasis at this time.

	\balance

\end{document}